%\pagebreak
\section{Atomic Compare}
\label{sec:cas}

\index{constructs!atomic@\kcode{atomic}}
\index{atomic construct@\kcode{atomic} construct}
\index{clauses!capture@\kcode{capture}}
\index{clauses!compare@\kcode{compare}}
\index{capture clause@\kcode{capture} clause}
\index{compare clause@\kcode{compare} clause}

The \kcode{compare} clause was added to the \plc{extended-atomic} clauses
for C/C++ in OpenMP 5.1 and extended to Fortran in OpenMP 6.0.
The \kcode{compare} clause extends the semantics to perform the \kcode{atomic}
update conditionally. 

In the following example, two formats of structured blocks
are shown for associated \kcode{atomic} constructs with the \kcode{compare} clause.
In the first \kcode{atomic} construct, the format forms a \emph{conditional update} statement.
In the second \kcode{atomic} construct the format forms a \emph{conditional expression} statement.
The ``greater than'' and ``less than'' forms and 
the conditional expression form from Fortran 2023
are available to the \kcode{compare} clause for Fortran in OpenMP 6.0.
% One can use the \vcode{max} and \vcode{min} functions with the \kcode{atomic update}
% construct to perform the C/C++ example operations.

\cexample[5.1]{cas}{1}
\ffreeexample[6.0]{cas}{1}
\clearpage

In OpenMP 5.1 the \kcode{compare} clause with the \kcode{capture} clause
was added to support \emph{Compare And Swap} (CAS) semantics 
in C/C++ and extended to Fortran in OpenMP 6.0.
In the following example, the \ucode{enqueue} routine
(a naive implementation of a Michael and Scott enqueue function), uses the
\kcode{compare} with \kcode{capture} clause to perform compare
(\plc{x == e}) and swap (\plc{if-else} assignments) of the
form in C/C++: 
\begin{description}[noitemsep,labelindent=5mm,widest=f90]
\item \splc{{ r = x == e; if (r) { x = d; } else { v = x; } }}.
\end{description}

The equivalent form for Fortran pointers is:
\begin{description}[noitemsep,labelindent=5mm,widest=f90]
\item \splc{r = associated(x, e); if (r) then; x => d; else; v => x; end if}
\end{description}
where the intrinsic function \bcode{associated(\plc{x, e})} is used
for pointer comparison.
In the example code, the corresponding variables for ``\plc{x, e, d, v}''
are C/C++ variables ``\ucode{queue->tail}, \ucode{queue->next}, 
\ucode{node}, \ucode{node->next}'', 
or Fortran variables ``\ucode{queue\%tail}, \ucode{queue\%next}, 
\ucode{node}, \ucode{node\%next}'', respectively.
The example program concurrently enqueues nodes from an array of nodes
(\ucode{nodes[N]} or \ucode{nodes(N)}).

\cexample[5.1]{cas}{2}
\ffreeexample[6.0]{cas}{2}
